\usepackage{geometry}
\geometry{top=3cm, bottom=2cm, left=3cm, right=3cm}

% \usepackage{bookmark}

% Fonts (XeLaTeX)
\setCJKmainfont[AutoFakeBold=4]{標楷體}
\setCJKfamilyfont{kai}[AutoFakeBold=4]{標楷體}

% Paragraph spacing
\setlength{\parindent}{2em}
\setlength{\parskip}{1em}
\usepackage{indentfirst} % 确保首段缩进

% Headers/footers
\usepackage{fancyhdr}
\fancypagestyle{plain}{
  \fancyhf{}
  \fancyfoot[C]{\thepage}
  \renewcommand{\headrulewidth}{0pt}
}

% Section formatting
\usepackage{titlesec}
\newcommand{\myparagraph}[1]{\paragraph{#1}\mbox{}\\}
\newcommand{\mysubparagraph}[1]{\subparagraph{#1}\mbox{}\\}

\titleformat{\paragraph}
  {\normalfont\normalsize\bfseries}
  {第\chinese{paragraph}款}
  {1em}
  {}
\titlespacing*{\paragraph}{0em}{0em}{0em}

\titleformat{\subparagraph}
  {\normalfont\normalsize\bfseries}{第\chinese{subparagraph}目}{1em}{}
\titlespacing*{\subparagraph}{1em}{0em}{0em}

% 設定章節編號格式
\ctexset{
  section = {
    format = \centering\zihao{-2}\CJKfamily{kai},
    name = {第,章},
    number = \chinese{section},
    aftername = \quad
  },
  subsection = {
    format = \zihao{3}\CJKfamily{kai},
    name = {第,節},
    number = \chinese{subsection},
    aftername = \quad
  },
  subsubsection = {
    format = \zihao{-3}\CJKfamily{kai},
    name = {第,項},
    number = \chinese{subsubsection},
    aftername = \quad
  },
  paragraph = {
    format = \zihao{5}\CJKfamily{kai},
    name = {第,款},
    number = \chinese{paragraph},
    aftername = \quad,
    indent = 0pt
  },
  subparagraph = {
    format = \zihao{5}\CJKfamily{kai},
    name = {第,目},
    number = \chinese{subparagraph},
    aftername = \quad,
    indent = 0pt
  }
}
\setcounter{secnumdepth}{5}
\setcounter{tocdepth}{5}

% Tables
\usepackage{longtable}   % 支持跨页表格
\usepackage{tabularray}
\usepackage{threeparttable} % 讓表格內 footnote 正常顯示
\usepackage{siunitx}
\sisetup{detect-all}
\usepackage{arydshln}    % Provides the \hdashline command
\usepackage{array}       % For advanced column specifications
\usepackage{tabularx}
\newcolumntype{H}{>{\hspace{-9999pt}}c}
\newcolumntype{L}[1]{>{\raggedright\arraybackslash}p{#1}}
\newcolumntype{Y}{>{\raggedright\arraybackslash}X}
\usepackage{caption}     % Enhances caption formatting
\usepackage{booktabs} % Optional: For better-looking tables
\usepackage{graphicx} % Optional: For including images in tables or elsewhere
\usepackage{float}    % Optional: To control table placement
\usepackage{multirow} % Required for multirow functionality in tables
\usepackage{pdflscape}

% Lists
\usepackage{enumitem}
\setlist[itemize]{topsep=0em, partopsep=0pt, parsep=0pt, itemsep=0.5em}

% Bibliography
% Fall back to a stock biblatex style when the custom NTU Law style is not
% installed on the current machine, so the document remains compilable.
\newif\ifhasntulawstyle
\IfFileExists{ntulaw.bbx}{\hasntulawstyletrue}{\hasntulawstylefalse}

\ifhasntulawstyle
  \usepackage[backend=biber,style=ntulaw,citestyle=ntulaw]{biblatex}
\else
  \usepackage[backend=biber,style=authortitle,citestyle=authortitle]{biblatex}
\fi

\IfFileExists{My_Bib.bib}{\addbibresource{My_Bib.bib}}{}
\IfFileExists{1121財政法.bib}{\addbibresource{1121財政法.bib}}{}
\IfFileExists{Web_Sources.bib}{\addbibresource{Web_Sources.bib}}{}

\ifhasntulawstyle
  % Unified citation helpers for this thesis.
  % Use these in Chinese footnotes when chaining multiple foreign full citations,
  % so the note ends with one Chinese full stop instead of repeated English periods.
  \DeclareCiteCommand{\zhfullciteNP}
    {\usebibmacro{prenote}}
    {\begingroup\renewcommand*{\addperiod}{}\renewcommand*{\ntufinalpunct}{}\usedriver{}{\thefield{entrytype}}\endgroup}
    {\multicitedelim}
    {\usebibmacro{postnote}}
  \DeclareMultiCiteCommand{\zhfullcites}{\zhfullciteNP}{\printtext{；\space}}

  % Chinese article helper:
  % Use for Chinese journal articles when ntulaw's language branching is unreliable.
  \DeclareCiteCommand{\zharticleciteNP}
    {\usebibmacro{prenote}}
    {\begingroup\renewcommand*{\ntufinalpunct}{}\usebibmacro{article:chinese}\endgroup}
    {\multicitedelim}
    {\usebibmacro{postnote}}
  \DeclareMultiCiteCommand{\zharticlecites}{\zharticleciteNP}{\printtext{；\space}}
\else
  \newcommand{\zhfullciteNP}{\fullcite}
  \newcommand{\zhfullcites}{\fullcite}
  \newcommand{\zharticleciteNP}{\fullcite}
  \newcommand{\zharticlecites}{\fullcite}
\fi

% TOC and hyperlinks
\usepackage{shorttoc}
\usepackage{hyperref}
\usepackage{xurl}

% Boxes
\usepackage{tcolorbox}
\tcbuselibrary{breakable,skins}
\definecolor{notered}{HTML}{F2AFC0}
\definecolor{noteblue}{HTML}{A9C9F5}
\definecolor{notegreen}{HTML}{B9D9A8}
\definecolor{noteorange}{HTML}{F7E39C}
\definecolor{notepurple}{HTML}{D8B8E8}
\definecolor{noteteal}{HTML}{AEE1D6}
\definecolor{notegray}{HTML}{D7DADF}
\definecolor{critiquebg}{HTML}{F6D6DD}

\tcbset{
  noteboxbase/.style={
    enhanced,
    breakable,
    coltitle=black,
    fonttitle=\bfseries,
    boxrule=0.5mm,
    arc=3mm,
    left=2mm,
    right=2mm,
    top=1.5mm,
    bottom=1.5mm,
    boxsep=2mm,
    before skip=0.8em,
    after skip=0.8em,
    before upper={\setlength{\parskip}{0.6em}\setlength{\parindent}{0pt}},
    drop shadow={fill=black!12!white}
  }
}

\newtcolorbox{notebox}[2][]{
  noteboxbase,
  colframe=#1,
  colback=white,
  colbacktitle=#1!18!white,
  coltitle=black,
  title={#2}
}

% \newcommand{\noteplaceholder}{%
%   \textit{［此框內容尚未填入］}%
% }

\newcommand{\sep}{%
  \par\noindent
  {\color{black!35}\leavevmode\leaders\hbox{\rule{2.2pt}{0.4pt}\hskip1.6pt}\hfill\kern0pt\par}%
}

\newcommand{\notecolor}{notegray}

\newcommand{\callout}[1]{%
  \par\vspace{0.5em}%
  \noindent
  \colorbox{\notecolor!18!white}{%
    \parbox{\dimexpr\linewidth-2\fboxsep\relax}{%
      \setlength{\parindent}{0pt}%
      #1%
    }%
  }%
  \par\vspace{0.2em}%
}

\newcommand{\icallout}[1]{%
  \tcbox[
    on line,
    boxrule=0pt,
    arc=1pt,
    boxsep=0pt,
    left=1.5pt,
    right=1.5pt,
    top=1pt,
    bottom=1pt,
    colback=\notecolor!18!white,
    nobeforeafter,
    tcbox raise base
  ]{#1}%
}

\newcommand{\inlinecallout}[1]{\icallout{#1}}

\newcommand{\critique}[1]{%
  \par\vspace{0.5em}%
  \noindent
  \colorbox{critiquebg}{%
    \parbox{\dimexpr\linewidth-2\fboxsep\relax}{%
      \setlength{\parindent}{0pt}%
      #1%
    }%
  }%
  \par\vspace{0.2em}%
}

\newcommand{\icritique}[1]{%
  \tcbox[
    on line,
    boxrule=0pt,
    arc=1pt,
    boxsep=0pt,
    left=1.5pt,
    right=1.5pt,
    top=1pt,
    bottom=1pt,
    colback=critiquebg,
    nobeforeafter,
    tcbox raise base
  ]{#1}%
}

\newcommand{\noteannotation}[2]{%
  \par\vspace{0.5em}%
  \noindent
  \colorbox{#1!18!white}{%
    \parbox{\dimexpr\linewidth-2\fboxsep\relax}{%
      \setlength{\parindent}{0pt}%
      \footnotesize #2%
    }%
  }%
}

\NewDocumentCommand{\noteboxcmd}{m O{} +m g}{%
  \def\notecolor{#1}%
  \if\relax\detokenize{#2}\relax
    \begin{notebox}[#1]{}%
      #3%
      \IfNoValueF{#4}{\noteannotation{#1}{#4}}%
    \end{notebox}%
  \else
    \begin{notebox}[#1]{#2}%
      #3%
      \IfNoValueF{#4}{\noteannotation{#1}{#4}}%
    \end{notebox}%
  \fi
}

\NewDocumentCommand{\emptynoteboxcmd}{m O{}}{%
  \if\relax\detokenize{#2}\relax
    \begin{notebox}[#1]{}%
      \noteplaceholder
    \end{notebox}%
  \else
    \begin{notebox}[#1]{#2}%
      \noteplaceholder
    \end{notebox}%
  \fi
}

\NewDocumentCommand{\rednote}{O{} +m g}{\noteboxcmd{notered}[#1]{#2}{#3}}
\NewDocumentCommand{\bluenote}{O{} +m g}{\noteboxcmd{noteblue}[#1]{#2}{#3}}
\NewDocumentCommand{\greennote}{O{} +m g}{\noteboxcmd{notegreen}[#1]{#2}{#3}}
\NewDocumentCommand{\orangenote}{O{} +m g}{\noteboxcmd{noteorange}[#1]{#2}{#3}}
\NewDocumentCommand{\purplenote}{O{} +m g}{\noteboxcmd{notepurple}[#1]{#2}{#3}}
\NewDocumentCommand{\tealnote}{O{} +m g}{\noteboxcmd{noteteal}[#1]{#2}{#3}}
\NewDocumentCommand{\graynote}{O{} +m g}{\noteboxcmd{notegray}[#1]{#2}{#3}}
\NewDocumentCommand{\note}{O{} +m g}{\noteboxcmd{notegray}[#1]{#2}{#3}}

\NewDocumentCommand{\emptyrednote}{O{}}{\emptynoteboxcmd{notered}[#1]}
\NewDocumentCommand{\emptybluenote}{O{}}{\emptynoteboxcmd{noteblue}[#1]}
\NewDocumentCommand{\emptygreennote}{O{}}{\emptynoteboxcmd{notegreen}[#1]}
\NewDocumentCommand{\emptyorangenote}{O{}}{\emptynoteboxcmd{noteorange}[#1]}
\NewDocumentCommand{\emptypurplenote}{O{}}{\emptynoteboxcmd{notepurple}[#1]}
\NewDocumentCommand{\emptytealnote}{O{}}{\emptynoteboxcmd{noteteal}[#1]}

% 設置 quote 框框樣式
\tcbset{
  myquote/.style={
    colframe=black,
    colback=gray!10,
    coltitle=black,
    fonttitle=\itshape,
    sharp corners,
    boxrule=0.5mm,
    width=\linewidth,
    leftrule=1mm,
    title={出處：某書籍或文件},
    before skip=1em,
    after skip=1em
  }
}

\renewcommand{\contentsname}{詳目}
\renewcommand{\listfigurename}{圖目錄}
\renewcommand{\listtablename}{表目錄}


\usepackage{makecell}

\newcommand{\celltop}[1]{%
  \parbox[t]{\linewidth}{%
    \raggedright
    \strut #1\strut
  }%
}
